\section{Elecciones finales: la eficiencia es una apariencia, las interfaces de estado son la esencia}

Los ejemplos anteriores de modelos, tokens y H-Net sirven para una corrección final: los SSM no son solo "más rápidos pero más débiles", y la atención no es solo "más fuerte pero cuadrática". Sus ventajas y desventajas surgen conjuntamente de la forma del autoregressive state. Las últimas cuatro figuras primero reiteran las resoluciones adecuadas para la atención, luego reescriben por separado el tradeoff entre SSM y atención, y finalmente comprimen la investigación arquitectónica en un problema de FLOPs a capacidades.

\subsection{Lo que realmente hace fuerte a la atención es la pre-compresión de abstracción}

Esta sección regresa al problema de abstracción que ha aparecido repetidamente en toda la clase. El cómic dibuja la vergüenza del Transformer sobre raw high-resolution input como un malentendido: el modelo no entiende nativamente ninguna granularidad; es más fuerte en unidades significativas, combinables y direccionables individualmente. El texto rojo a la derecha de la figura no niega la atención, sino que le asigna una ubicación: primero compresione los datos hasta el nivel adecuado y luego utilice el direccionamiento global del contenido.

\lecturefigure{slide-030.jpg}{Área más efectiva de la atención: datos pre-comprimidos y nivel de abstracción correcto}{Grabación oficial de Stanford, 00:49:20--00:50:09}

\begin{importantbox}{La selección arquitectónica debe ser consciente de la resolución}
Una misma modalidad puede utilizar diferentes primitivos en distintas capas: las capas de píxeles/bytes realizan primero compresión local o recurrente, las capas de objetos/palabras/tramos realizan luego atención, y finalmente las capas de decisión pueden conectarse a memoria externa. Forzar todo el pipeline a ser un único tipo de capa suele estar haciendo que la arquitectura se adapte a unidades incorrectas.
\end{importantbox}

\subsection{SSM: tache "eficiencia", conserve statefulness y compresión}

En la figura, primero se escribe "escalado lineal tiempo, inferencia tiempo constante", luego se tacha efficiency y se reemplaza por model stateful / compressive, etiquetando brain-like state. El orador no dice que la eficiencia no sea importante, sino que es insuficiente para explicar por qué este tipo de modelos pueden ser más adecuados en state tracking, interacción online y construcción de abstracción; estas capacidades y las debilidades de retrieval provienen del estado finito.

\lecturefigure{slide-031.jpg}{Elección final SSM: compresión stateful, tracking de estado y pérdida de retrieval granular}{Grabación oficial de Stanford, 00:50:10--00:51:19}

\teachervoice{00:50:05--00:51:27, Gu dice que efficiency es un red herring: los pros y contras de los SSM son dos caras de la misma moneda, es decir, brain-like autoregressive state. El estado finito pierde detalles precisos y crea presión para el seguimiento y la abstracción.}

\subsection{Atención: tache "cuadrático", conserve dependencia de resolución}

La fortaleza de la atención es el acceso granular al context pasado, especialmente adecuado para recall y retrieval. En la figura se tacha "escalado cuadrático", no para cancelar la complejidad, sino para señalar una limitación más profunda: el modelo es altamente sensible a la resolución semántica de los tokens de entrada; un cache tipo database guardaría fielmente las unidades que se le dan, pero no decidiría por el diseñador si esas unidades valen la pena ser guardadas.

\lecturefigure{slide-032.jpg}{Elección final atención: cache tipo database, acceso preciso y dependencia token-resolución}{Grabación oficial de Stanford, 00:51:20--00:53:09}

\begin{warningbox}{La atención cuadrática no es "desperdicio" cuando se requiere memoria precisa}
Si la tarea realmente exige conservar ítem por ítem cláusulas contractuales, tokens de código, párrafos de evidencia o retornos de herramientas externas, un cache $O(L)$ y lectura por paso $O(L)$ pueden ser el costo informativo necesario para completar la tarea. La pregunta correcta no es "¿podemos eliminar lo cuadrático?", sino "qué capas, qué tokens y qué escalas temporales merecen pagar este costo".
\end{warningbox}

\begin{table}[H]
\centering
\caption{Elección simétrica entre dos tipos de autoregressive state}
\begin{tabularx}{\textwidth}{L{0.18\textwidth}>{\raggedright\arraybackslash}X >{\raggedright\arraybackslash}X}
\toprule
Eje de comparación & Atención token-resolución & Estado recurrente comprimido \\
\midrule
Estado entre pasos & KV cache que crece con $L$ & Estado finito desacoplado de $L$ \\
Acceso preciso & Fuerte; posiciones históricas direccionables individualmente & Débil; detalles mezclados en estado compartido \\
Actualización online & Lectura de cache histórico por paso & Actualización recursiva mediante interfaz fija \\
Inducción sesgo & database, retrieval, copy & tracking, smoothing, compresión \\
Factores sensibles & resolución token y unidad semántica & capacidad estado y expresividad actualización \\
Remedio típico & MQA/GQA, compresión KV, enrutamiento retrieval & Ampliar estado, selectividad, inserción atención \\
\bottomrule
\end{tabularx}
\end{table}

\subsection{Pregunta central de la investigación arquitectónica: ¿cada FLOP se gasta en estructura significativa?}

La última figura dibuja el sistema de entrenamiento como una caja negra: a la izquierda se inyectan FLOPs y datos, a la derecha se obtienen capacidades. El objetivo real no es minimizar FLOPs ni maximizar el modelo, sino mejorar la eficiencia de conversión de cómputo a capacidades. Si el Transformer cacheara cada carácter en un flujo de caracteres pero rara vez necesitara direccionamiento de carácter individual, esos FLOPs no coinciden con la estructura de la tarea; si la tarea requiere recordar con precisión cada detalle, pagar el costo de atención sería razonable.

\lecturefigure{slide-033.jpg}{FLOPs $\rightarrow$ modelo $\rightarrow$ capacidades: ¿usa el modelo sabiamente el cómputo?}{Grabación oficial de Stanford, 00:53:10--00:54:34}

\[
\text{system efficiency}
=\frac{\text{useful capabilities}}{\text{training dollars}}
=\frac{\text{FLOPs}}{\text{dollars}}
\times
\frac{\text{capabilities}}{\text{FLOPs}}.
\]
El primer término está influenciado por hardware, kernel, paralelismo y utilización; el segundo término está influenciado por arquitectura, datos, objetivo y representación. Gu estudia principalmente el segundo término, pero los sistemas de despliegue deben optimizar ambos simultáneamente. Un FLOP es matemáticamente idéntico, pero estructuralmente puede ser completamente diferente: puede usarse para comparar chunks significativos o para procesar repetidamente unidades de alta resolución sin sentido.

\begin{importantbox}{Principio de diseño final}
No seleccione primero la arquitectura y fuerce a los datos a adaptarse, sino dibuje primero la estructura de la información en diferentes escalas temporales: qué detalles deben conservarse, cuáles pueden comprimirse, cuáles requieren acceso aleatorio y cuáles solo necesitan seguimiento de estado. Luego coloque atención, SSM, convolución, codificador local y memoria externa en las capas más coincidentes.
\end{importantbox}

\subsection{Q\&A: backprop, hardware y "óptimo local"}

Un estudiante cuestiona que los SSM modernos sacrifican expresividad recurrente para permitir backprop-through-time y paralelismo GPU, mientras que el cerebro humano no guarda muchas copias para realizar propagación hacia atrás. La respuesta de Gu es cautelosa: cree que la arquitectura actual, backprop y hardware podrían estar en un óptimo local común; la memoria de activación del backprop imponería una limitación física a la longitud de dependencias aprendibles, pero no está seguro si las tareas reales ya pueden sortear esto mediante medios de ingeniería.

\teachervoice{00:55:10--00:57:22, Gu cree que replantear el algoritmo de aprendizaje y la restricción de hardware podría producir modelos fundamentalmente mejores; al mismo tiempo admite no haber ofrecido una solución madura alternativa al backprop.}

\begin{warningbox}{Preguntas abiertas no son defectos confirmados}
El hecho de que "el cerebro humano no use backprop" no implica que el backprop sea inútil para el aprendizaje automático; el hecho de que "la memoria no quepa secuencias ultra largas" tampoco significa que los modelos no puedan aprender absolutamente dependencias largas, ya que currículum, tareas sintéticas, retrieval y entrenamiento recurrente pueden ofrecer rutas indirectas. Las ideas de la clase deben servir como hipótesis de investigación, no como conclusiones.
\end{warningbox}

\subsection{Q\&A: límites de H-Net, memoria y modelos diminutos}

Sobre H-Net, el orador afirma claramente que la versión actual es solo un primer paso funcional, sin validación en la escala máxima. Los chunks dinámicos no tienen vocabulario infinito: los bytes iniciales entran en embeddings limitados, mientras que lo subsiguiente selecciona boundaries y representa chunks agrupados en el espacio latente. Para contextos largos, imagina una compresión recursiva hacia niveles de memoria cada vez más altos, permitiendo que la atención del modelo interno se enfoque en entradas menos numerosas pero más significativas, aunque aún no hay validación posterior.

\teachervoice{00:59:09--01:03:31, la representación de chunk no pasa por lookup de vocabulario, sino que preserva posiciones seleccionadas o resúmenes agrupados en el espacio latente del codificador; la memoria jerárquica es una dirección futura coherente con la filosofía de H-Net, no un resultado actual.}

Para data-constrained / modelos diminutos, Gu responde directamente "no sé cuál es la regla general". Solo menciona que un pequeño experimento Mamba-3 realizado en ese momento aún descubrió que poca atención es muy importante. Esta incertidumbre en sí misma merece ser conservada: los inductivos sesgos pueden cambiar en diferentes regímenes de escalado y no se puede retroinferir desde una curva de mil millones de parámetros a un modelo diminuto.

\subsection{Q\&A: ¿pueden los SSM explicarse como un mapa de atención?}

Algunas formas de SSM / atención lineal pueden expandirse en un mapa de dependencia token-a-token, visualmente similar a un mapa de atención. La observación de Gu es que la atención softmax a veces se manifiesta como una atención aproximadamente dura, concentrándose en pocos tokens; la dependencia del SSM suele ser más difusa, coincidiendo con la intuición de "mezclar información dentro del estado". Pero esto es solo un proxy; la interpretabilidad mecánica aún está en curso.

\teachervoice{01:04:52--01:06:29, el orador dice que la visualización tipo atención puede revelar diferencias de comportamiento entre dos tipos de modelos, pero no es la única herramienta de explicación; la interpretabilidad mecánica completa del SSM sigue siendo trabajo abierto.}

\subsection{Resumen de la sección}

La comparación final ya no es "modelo cuadrático versus modelo lineal", sino cache tipo database versus estado comprimido tipo cerebro. La atención cambia acceso preciso a cambio de costo por token, el SSM cambia estado online y presión de abstracción a cambio de compresión irreversible; híbridos y jerarquías permiten que ambos se especialicen por resolución. Lo que la investigación arquitectónica debe optimizar al final es el grado de coincidencia entre cada FLOP y la estructura informativa de la tarea.
